\documentclass{article}
\newcommand{\bpnode}[1]{}
\newtheorem{definition}{Definition}
\newtheorem{theorem}{Theorem}
\newtheorem{proof}{Proof}
\title{Continuous maps and their composition}
\author{Qprint demonstration notes}
\begin{document}
\section{Continuous maps}
\subsection{The identity map}
\bpnode{identity}
\begin{definition}
Let $X$ be a set. The \emph{identity map} on $X$ is the function
\[
  \operatorname{id}_X : X \longrightarrow X, \qquad x \longmapsto x.
\]
It fixes every point of the underlying space.
\end{definition}

The same construction is available in Lean, Agda and Coq. No topological
structure is needed to define this map.

\bpnode{continuous-id}
\begin{theorem}
Let $X$ be a topological space. The identity map
$\operatorname{id}_X : X \to X$ is continuous.
\end{theorem}

Indeed, for every open subset $U \subseteq X$, we have
\[
  \operatorname{id}_X^{-1}(U) = U.
\]
Since $U$ is open, the criterion for continuity is satisfied.

\subsection{Composition}
\bpnode{continuous-comp}
\begin{theorem}
Let $X$, $Y$ and $Z$ be topological spaces, and let
$f : X \to Y$ and $g : Y \to Z$ be continuous maps.
Then their composition
\[
  g \circ f : X \longrightarrow Z
\]
is continuous.
\end{theorem}

This closure property makes it possible to build complicated continuous
maps from simpler ones. In the formal statement, continuity of both
maps appears as an explicit hypothesis.

\bpnode{composition-proof}
\begin{proof}
Take an open subset $W \subseteq Z$. By continuity of $g$,
the preimage $g^{-1}(W)$ is open in $Y$.
Continuity of $f$ then implies that $f^{-1}(g^{-1}(W))$ is open in $X$.
Using the identity
\[
  (g \circ f)^{-1}(W) = f^{-1}(g^{-1}(W)),
\]
we conclude that $g \circ f$ is continuous.
\end{proof}

\section{A categorical perspective}
\bpnode{top-category}
\begin{theorem}
Topological spaces and continuous maps form a category.
\end{theorem}

The preceding results supply the identity morphisms and closure under
composition. The remaining laws follow from the corresponding laws
for functions:
\begin{itemize}
\item Composition of functions is associative.
\item The identity map is a left and right unit for composition.
\end{itemize}
\end{document}
